% sections/findings.tex
\section{C3: What moved the score, priced in steps}
\label{sec:findings}

More optimizer steps in the same 30 minutes (\cref{sec:steps}) and a handful of structural changes drove most of the
descent in \cref{fig:history}. This section prices the changes with $\kappa$ (\cref{sec:price}), gives each one a
noise scale (\cref{sec:levers}), examines the row pool, attention-source reuse and the EMA prewarm
(\cref{sec:rowpool,sec:asr,sec:prewarm}) and the noise and selection effects (\cref{sec:seeds,sec:salts}), and ends
with what did not work (\cref{sec:negative}) and what the remaining gap would cost (\cref{sec:gap}).

\subsection{Levers, their noise and their price}
\label{sec:levers}

\Cref{tab:levers} lists the changes that helped. The noise column is the SD of a single measured difference under the
order-redraw noise model of \cref{sec:seeds}: $\sqrt2$ times the shuffle-salt SD for a chip pair
(\val{noise-chip-pair}), plus the offset noise of two uploads for an official pair (\val{noise-official-pair}),
divided by $\sqrt n$ for $n$ pairs. It is the right scale for changes that act after the warm-up; a change that
re-rolls the warm-up path faces the larger seed noise. Four of the six largest effects (multi-token prediction,
batch warm-up, LeakyReLU$^2$, the key offset) come from the campaign log and have no run identifiers in the
repository; the late levers are single pairs at seed~73. Only the row pool, the
attention-source reuse and the earlier, larger levers clear about three noise SDs.

\begin{table}[tbp]
  \centering
  \caption{\textbf{Levers that worked.} $\Delta$: change in \valbpb{} against the lever's control, in $10^{-3}$
    (negative is better). Noise SD: SD of the measured difference ($10^{-3}$, see text). Evidence: ``chip pair'',
    same-seed same-chip rehearsals, raw; ``step-norm.'', normalised for a step-count difference; ``official pair'',
    two uploads that differ by this change. Steps: the step equivalent $\exp(|\Delta|/\kappa)-1$ at
    $\kappa=\val{kappa}$. \textsuperscript{$\dagger$}\,Listed in \texttt{docs/FINDINGS.md} without run identifiers,
    mostly on earlier recipe bases. \textsuperscript{$\ddagger$}\,The batch warm-up changes the step count itself, so
    a step equivalent would double-count.}
  \label{tab:levers}
  % GENERATED by paper/analysis -- levers that worked, docs/FINDINGS.md sections 1 and 3; step equivalent at kappa 0.057
{\footnotesize\setlength{\tabcolsep}{4pt}
\begin{tabular}{@{}p{4.2cm}rrrlrrp{3.3cm}@{}}
\toprule
Lever & $\Delta$ & Noise SD & $|\Delta|/\text{SD}$ & Evidence & $n$ & Steps & Note \\
\midrule
Multi-token prediction (3 tokens, faded) & \ensuremath{-}4.3 & 0.49 & 8.7 & campaign log$^{\dagger}$ & -- & 7.8\% &  \\
Row pool, 256 micro-batches & \ensuremath{-}2.3 & 0.49 & 4.6 & chip pair & 1 & 4.1\% & chip E; official $-$2.2; $-$0.5\% steps \\
Batch warm-up 65k$\to$131k$\to$262k & \ensuremath{-}1.9 & 0.49 & 3.8 & campaign log$^{\dagger}$ & -- & n/a & changes steps itself$^{\ddagger}$ \\
LeakyReLU(0.35)$^2$ MLP & \ensuremath{-}1.7 & 0.49 & 3.4 & campaign log$^{\dagger}$ & -- & 3.0\% & chip C vs ReLU$^2$ $-$1.05; official $-$1.5 \\
Attention-source reuse (6--8 reuse 5) & \ensuremath{-}1.5 & 0.25 & 6.1 & chip pair & 4 & 2.7\% & 3 seeds, 2 chips; 5th pair (s97) $+$1.15 raw; step cost $<$0.1\% \\
Rotary key offset & \ensuremath{-}1.3 & 0.49 & 2.6 & campaign log$^{\dagger}$ & -- & 2.3\% &  \\
Cooldown 0.7 + small-batch LR & \ensuremath{-}0.6 & 0.29 & 2.1 & chip, step-norm. & 3 & 1.1\% & all 3 negative; seeds differ \\
EMA blend 0.6 (no prewarm) & \ensuremath{-}0.5 & 0.49 & 1.0 & chip pair & 1 & 0.9\% & official $-$0.3 (K77a$-$K73s4) \\
MLP output-projection LR $\times$0.8 & \ensuremath{-}0.4 & 0.49 & 0.8 & chip, step-norm. & 1 & 0.7\% & seed 73 only \\
EMA prewarm (K82s4 vs K77a) & \ensuremath{-}0.3 & 0.73 & 0.4 & official pair & 1 & 0.5\% & recovers 15--30 steps \\
\bottomrule
\end{tabular}
}

\end{table}

Three entries need a word. \emph{LeakyReLU(0.35)$^2$} appears with three sizes because they are three comparisons:
\val{lever-leaky}\docnum{} in the campaign log on an earlier base, \val{leaky-chip-c}\docnum{} in a chip-C pair against
ReLU$^2$ at equal steps (the comparison the GPU proxy reproduces), and \val{leaky-official} officially
(\upl{K59}$-$\upl{K57}). \emph{Cooldown~0.7} is reported here together with the lower small-batch learning rate, as three
step-normalised chip pairs\docnum{}, all three negative; \cref{tab:gpu} uses two raw chip pairs of the cooldown change alone,
of opposite sign. The official \upl{K63}$-$\upl{K60} delta (\val{cooldown-official}) is confounded, because
\upl{K63} also turned the EMA off and added 2~s of budget. \emph{Official deltas} between uploads published to four
decimals carry a rounding error of up to $\pm2\times$\val{official-rounding}.

\subsection{The row pool: decorrelating a streaming loader}
\label{sec:rowpool}

\paragraph{Mechanism.} The kit's loader fills each row of a micro-batch from one of several parallel token streams,
so the rows of consecutive micro-batches are contiguous slices of the same few streams. From
micro-batch~\val{rowpool-from} on, our loader instead draws rows without replacement from a look-ahead pool of
\val{rowpool-size} micro-batches (\knob{FF_ROW_SHUFFLE}), a shuffle buffer at row granularity
\citep[cf.][]{xu2022corgipile,mishchenko2020reshuffling}. The first \val{rowpool-from} micro-batches, which cover the
warm-up, are unchanged, so that the seed's warm-up path (\cref{sec:seeds}) is not re-rolled. The shuffle has its own
RNG with a salt, which re-draws the order after micro-batch~\val{rowpool-from} without touching initialisation.

\paragraph{Evidence.} On chip E, the pool scored \val{rowpool-chip-treat}\docnum{} at \val{rowpool-steps-treat}\docnum{}
steps against \val{rowpool-chip-ctrl}\docnum{} at \val{rowpool-steps-ctrl}\docnum{} for the same file without it:
\val{rowpool-chip}\docnum{} raw, about \val{rowpool-stepnorm}\docnum{} after normalising for the
\val{rowpool-step-cost}\docnum{} of steps the pool cost, and about \val{rowpool-z} noise SDs. Officially, \upl{K70a}
scored \val{official-k70a} against \upl{K63}'s \val{official-k63}, a difference of \val{rowpool-official}. That
agreement is what determinism predicts and is not independent evidence that the effect generalises: both numbers come
from the same trajectory pair at seed~73. And because the pool changes the data order, its measured effect contains
one draw of the order noise (SD about \val{salt-sd}\docnum{}). A pool of \val{pool-large}\docnum{} gave the same
result as \val{rowpool-size}; a pool of \val{pool-small}\docnum{} helped less. Other order manipulations were null
within $\pm\val{rowpool-null}$\docnum{}: larger pools, epoch reshuffles, pool-aware accumulation and document
reordering.

\paragraph{What we do not know.} The evidence is one chip pair at one seed and one salt. It is probably real at about
\val{rowpool-z} order-noise SDs, but it is not replicated, and the pair is not in the released run data. Our working hypothesis is
that the pool reduces within-batch correlation between rows of the same stream, and with it gradient variance, in the
spirit of shuffling results for SGD \citep{gurbuzbalaban2019reshuffling}. We have \emph{not} measured this; the checks
are listed in \cref{sec:limits}.

\subsection{Attention-source reuse}
\label{sec:asr}

Blocks 6, 7 and 8 reuse the attention source computed by block~5 (\knob{FF_ATTN_SRC=5:6,7,8}), in the spirit of
cross-layer attention sharing \citep{brandon2024cla,xiao2019sharing}. It is our best-replicated structural effect:
four same-seed chip pairs (seeds 73, 67 and 58 on chip~C, seed~73 on chip~D) gave \val{asr-pairs}\docnum{} (mean
\val{asr-mean}), at a step cost under \val{asr-step-cost}\docnum{}; a fifth pair, seed~97 on chip~D, was
\val{asr-s97-delta}\docnum{} raw, mostly a \val{asr-s97-step-deficit}\docnum{}-step deficit (\cref{tab:postfit}).
Officially, \upl{K60} improved on \upl{K59} by \val{asr-official}. The surrogate, whose coefficient for this feature
then rested on two runs, predicted the right sign on four of five post-fit attention-source pairs but overestimated
the size by about \val{asr-overestimate}\docnum{} (\cref{tab:postfit}).

\subsection{The EMA blend and the prewarm accounting trap}
\label{sec:prewarm}

\paragraph{The trap.} Our final weights are a 0.6 blend of the live weights with an exponential moving average
\citep{polyak1992averaging,moralesbrotons2024ema}. The EMA update is a separate graph. On a fresh instance, the first
EMA update compiles that graph \emph{inside} the charged clock and costs \val{ema-compile-steps}\docnum{} steps. A warm
rehearsal, run after an earlier EMA run on the same instance, finds the graph in the compiler cache and does not pay.
Warm rehearsals therefore looked \val{ema-warm-bias}\docnum{} bpb better than the cold official run would be.

\paragraph{A regression, not a discovery.} The fix, \knob{FF_EMA_PREWARM}, runs one EMA update during the excluded
start-up. It was not new on the final day: the ledger lists the EMA prewarm already in \upl{\val{first-prewarm}}
(26~September). It disappeared when \upl{\val{prewarm-dropped}} turned the EMA off, and \upl{K77a} re-enabled the EMA
without it. Only cold rehearsals exposed the regression, which is the general lesson below.

\paragraph{Evidence.} Step counts on one instance show the cost (\cref{tab:prewarm}). On chip~G the recipe without EMA
ran \val{steps-k73s6} steps, with EMA but without prewarm \val{steps-k77a}, and with both \val{steps-k82s7}; on chip~H,
without EMA \val{steps-k73s4} and with both \val{steps-k82s4}. These are single runs, and step counts jitter by
\val{jitter-range-lo}--\val{jitter-range-hi} steps within a seed sweep (SD \val{jitter-sd-lo}--\val{jitter-sd-hi};
\cref{sec:determinism}). The chip-G drop of \val{prewarm-g-drop} steps without prewarm is about \val{prewarm-g-drop-sds}
single-run jitter SDs but only just above the largest within-sweep range (\val{jitter-range-hi} steps), while the
difference between prewarm and no EMA (\val{prewarm-g-gain} steps on chip G, \val{prewarm-h-gain} on chip H) is inside
it. The salts differ between some of these files, which changes the data order but not the step time.
Officially, \upl{K82s4} (prewarm) scored \val{prewarm-official} against \upl{K77a} (no prewarm): one pair, below the
oracle's resolution (\cref{sec:resolution}), consistent with the step-based estimate. The EMA blend itself is a separate
effect of about \val{ema-blend-chip}\docnum{} in one chip pair without prewarm; officially \upl{K77a} improved on
\upl{K73s4} by \val{ema-blend-official}.

\begin{table}[tbp]
  \centering
  \caption{\textbf{Step counts around the compile trap} (single cold rehearsals, from
    \texttt{results/official-scores.csv}). Compare within a chip only; within-chip jitter is
    \val{jitter-range-lo}--\val{jitter-range-hi} steps.}
  \label{tab:prewarm}
  \small
  \begin{tabular}{llccr}
    \toprule
    Chip & Upload & EMA blend & EMA prewarm & Steps \\
    \midrule
    \val{chip-k73s6} & \upl{K73s6} & -- & -- & \val{steps-k73s6} \\
    \val{chip-k77a} & \upl{K77a} & \checkmark & -- & \val{steps-k77a} \\
    \val{chip-k82s7} & \upl{K82s7} & \checkmark & \checkmark & \val{steps-k82s7} \\
    \midrule
    \val{chip-k73s4} & \upl{K73s4} & -- & -- & \val{steps-k73s4} \\
    \val{chip-k82s4} & \upl{K82s4} & \checkmark & \checkmark & \val{steps-k82s4} \\
    \bottomrule
  \end{tabular}
\end{table}

\paragraph{The general lesson.} Any lazily compiled path that first runs after the clock starts is a hidden cost: EMA
updates, evaluation-mode switches, a late schedule phase with a new batch shape. On a compiled accelerator the compile
cost can exceed the benefit of the feature that triggers it, and only a \emph{cold} rehearsal reveals it. We recommend
that speed-run harnesses either compile all graphs before the clock starts or report compile time separately.

\subsection{Seeds and the warm-up coin flip}
\label{sec:seeds}

The largest single source of variance in our data was not a setting. In warm-up steps 3--10, some seeds take a
\emph{spike path}: the loss falls faster to step~4 and then much more slowly, ending step~10 roughly
\val{spike-gap-nats} nats higher than the good path (\cref{fig:warmup}a). The logged training loss is the weighted
multi-token-prediction sum (weights 1/0.5/0.25), not per-token cross-entropy, which is why it is near 13~nats. We
classify a run as spike-path when its step-10 loss is at least \val{spike-l10} nats. Good-path runs of these recipes
reach \val{l10-good} nats at step~10 and spike-path runs \val{l10-spike}; we chose the threshold after seeing the data,
in the empty gap between them, and one run at \val{l10-mid} is reported as intermediate. On the three clean seed sweeps
of a fixed recipe in our data (\cref{tab:seeds}), \val{spike-k} of \val{spike-n} seeds with a logged curve took the
spike path, \val{spike-frac} (Wilson 95\% interval \val{spike-wilson-lo}--\val{spike-wilson-hi}). The sweeps share
\val{spike-shared-seeds} seed, and the path may be a property of the seed, recipe and hardware together, so the runs
are not fully independent; counting each of the \val{spike-distinct-seeds} distinct seeds once gives
\val{spike-seed-k} spike seeds (\val{spike-seed-wilson}). Spike-path runs ended \val{spike-delta} worse on average
(95\% CI \val{spike-delta-lo} to \val{spike-delta-hi}), smaller than the \val{spike-penalty-doc}\docnum{} we had quoted
during the campaign from a broader, less controlled set of about \val{harvested-runs-doc}\docnum{} runs. We consider the
controlled estimate the better one.

\begin{figure}[tbp]
  \centering
  \includegraphics{warmup_paths.pdf}
  \caption{\textbf{The warm-up coin flip.} (a) Training loss (the weighted multi-token sum) over the first 20 steps
    for the \val{g7-nseeds} seeds of one recipe on one chip: spike-path seeds (dashed, triangles) drop faster to
    step~4, then stall. (b) Final score by path in the three seed sweeps; hollow points have no logged curve.}
  \label{fig:warmup}
\end{figure}

\paragraph{The variance budget.} Seeds re-draw the initialisation, the warm-up path and the data order. Within a
sweep, the seed SD is \val{seed-sd-all-lo}--\val{seed-sd-all-hi} including spike paths and
\val{seed-sd-good-lo}--\val{seed-sd-good-hi} among good-path seeds only. Shuffle salts re-draw only the order after
micro-batch~\val{rowpool-from}, with an SD of about \val{salt-sd}\docnum{}. Both are as large as most late
improvements (\cref{tab:levers}).

\paragraph{Selection.} Our seed, 73, was itself chosen by a lottery. In the 13-seed sweep of the \upl{K57} recipe it
ranks \val{g7-s73-rank} of \val{g7-nseeds}, \val{g7-s73-gap} better than the good-path mean, and every later
same-seed comparison is measured on that lottery winner. Effects measured on a selected seed need not hold at other
seeds \citep{cawley2010overfitting}. We followed two rules: a lever must act after step~19, so that it cannot re-roll
the warm-up path, and an effect under 0.0003 needs three seeds; the second was often not affordable. Our one official
test of a change to the first steps illustrates the first rule: an attention-input shift in blocks 6--8 (\upl{K65a})
scored \val{xshift-official} against its base. It may have re-rolled the warm-up path, but we did not log its
step-10 loss, and \val{xshift-official} is about \val{xshift-spike-ratio} times the controlled spike-path penalty
(\val{spike-delta}), so a genuinely harmful change cannot be ruled out.

\subsection{Salt lotteries and the winner's curse}
\label{sec:salts}

Because rehearsals are deterministic, a lottery over shuffle salts can be run on chip and the best rehearsal uploaded.
Four salts of the \upl{K73} recipe spanned \val{salt-span}\docnum{} on chip, and salt~4 was uploaded as \upl{K73s4};
the later uploads \upl{K77a} and \upl{K82s4} inherit it. A second lottery on the final recipe picked salt~7
(\upl{K82s7}); its size is not recorded in the released files. Where a lottery's draws ran on different chips, salt
and chip speed are confounded. Selection was only partly rewarded. Salt~4 kept its edge officially. Salt~6
rehearsed \val{salt6-reh-vs-salt4} against it, but on chip~G against chip~H, so that edge is confounded with the chip;
officially it scored \val{salt6-vs-salt4} against salt~4 (\upl{K73s6}). Salt~7 gave the best rehearsal of the campaign
(\val{salt7-rehearsal}) and tied salt~4 officially (\val{salt7-vs-salt4}).

A simple model explains the pattern. A salt's rehearsal advantage is its true effect (SD about \val{salt-sd-mid}) plus a
part that does not transfer to the official run, with SD $s_\eta$. The expected share of a selected advantage that
survives is the reliability ratio $\lambda$ (\cref{app:derivations}). With $s_\eta$ equal to the final-day offset SD,
$\lambda\approx\val{salt-shrink}$; that SD includes chip-to-chip variation, so we also use the within-chip offset SD
($\lambda=\val{salt-shrink-within}$), the since-\upl{K50} SD and a smaller 0.0002; together they give
$\lambda$ between \val{salt-shrink-lo} and \val{salt-shrink-hi}.
Lotteries are worth running when draws are cheap, but the projected gain should be shrunk by roughly a quarter to a half,
and if the official shard overlaps the public text (\cref{sec:contest}) part of what survives is selection on the
evaluation set itself \citep{dwork2015reusable,blum2015ladder}.

\subsection{What did not work}
\label{sec:negative}

\Cref{tab:negative} lists the changes that lost. Two patterns stand out. First, \emph{capacity changes lost in both
directions}: depth~10 and two key/value heads lost while running \val{depth10-steps}\docnum{} fewer steps, and depth~8
lost despite \val{depth8-steps}\docnum{} more steps. We read this as the 9-block model being a local optimum of the
capacity-for-steps trade at our kernels' speed, but it rests on one or two chip pairs without an equal-step comparison,
so we cannot say whether depth~10 lost \emph{only} because of its steps. Second, \emph{scalar retuning did not pay}:
none of \val{scalar-keys}\docnum{} scalar knob keys (learning-rate scales, betas, weight decay, initialisation,
momentum, soft cap, RoPE base, cooldown $\pm0.1$) won a confirmation. The four official tests on the final day
(\upl{K65a}--\upl{K65d}) were all worse than or tied with their base.

\begin{table}[tbp]
  \centering
  \caption{\textbf{What lost.} $\Delta$ against the control (positive is worse); evidence as in \cref{tab:levers}.
    Official deltas are against \upl{K63}. \textsuperscript{\S}\,May have re-rolled the warm-up path; its step-10 loss
    was not logged (\cref{sec:seeds}). Source: \texttt{docs/FINDINGS.md} \S5\docnum{} and
    \texttt{results/official-scores.csv}.}
  \label{tab:negative}
  % GENERATED by paper/analysis -- what lost, docs/FINDINGS.md section 5 and results/official-scores.csv (K65a-d)
{\small
\begin{tabular}{@{}p{7.6cm}rlr@{}}
\toprule
Change & $\Delta$ ($10^{-3}$) & Evidence & $n$ \\
\midrule
Depth 10 (bigger; 6--11\% fewer steps) & \ensuremath{+}3.8 to \ensuremath{+}4.3 & chip pair & 2 \\
Two KV heads (bigger) & \ensuremath{+}5.7 & chip pair & 1 \\
Depth 8 (smaller; +10.6\% steps) & \ensuremath{+}3.0 to \ensuremath{+}5.0 & chip pair & 2 \\
Late $k{=}8$ batch in the cooldown (K65b) & \ensuremath{+}0.5 & official pair & 1 \\
Split AdamW cooldown 0.4 (K65d) & \ensuremath{+}0.6 & official pair & 1 \\
EMA-Nesterov lookahead 0.3 (K65c) & \ensuremath{+}0.04 & official pair & 1 \\
AdamW LR $-$16\% & \ensuremath{+}0.8 to \ensuremath{+}0.9 & chip pair & 2 \\
AdamW warm-up reshaping (steps 0--19) & \ensuremath{+}1.0 & chip pair & 1 \\
Drop the block-0 MLP & \ensuremath{+}1.9 & chip pair & 1 \\
No positional encoding in layers 3--4 & \ensuremath{+}3.3 & chip pair & 1 \\
Attention-input shift, blocks 6--8 (K65a)$^{\S}$ & \ensuremath{+}3.4 & official pair & 1 \\
Scalar retunes (LR scales, betas, WD, init, momentum, soft cap, RoPE base, cooldown $\pm$0.1) & 0 of 61 keys won & chip pairs & 61 \\
Bigger pools, epoch reshuffles, pool-aware accumulation, document reordering & 0 $\pm$ 0.3 & chip pairs & -- \\
\bottomrule
\end{tabular}
}

\end{table}

\subsection{Pricing the gap in compute}
\label{sec:gap}

\Cref{tab:gap} converts the distance from our best official score, \val{best-official}, to reference points on the final
public leaderboard snapshot (about \val{leaderboard-time}\docnum{}; ranks only) into effective compute through
\cref{eq:exchange}. Closing the \val{gap-ten} gap to \#10 would have needed about \val{gap-ten-range} more effective
compute across our two calibrations of $\kappa$. The gap to \#1, \val{gap-one-range}, extrapolates far beyond the
\val{xr-anchor-pct} anchor and is an order of magnitude only.

\begin{table}[tbp]
  \centering
  \caption{\textbf{What each leaderboard reference would cost in effective compute}, from our best official score
    (\val{best-official}, \upl{K82s4}), as $\exp(\text{gap}/\kappa)-1$ for the price used ($\kappa=\val{kappa}$) and
    for the long-run calibration. Leaderboard values are from the final public snapshot, referred to by rank.}
  \label{tab:gap}
  % GENERATED by paper/analysis -- compute needed = exp(gap / kappa) - 1 (exchange.py)
{\small
\begin{tabular}{@{}lrrrr@{}}
\toprule
Target & val\_bpb & Gap & Compute, $\kappa=0.057$ & Compute, $\kappa=0.063$ \\
\midrule
\#10 on the final snapshot & 0.9554 & 0.0060 & +11\% & +10\% \\
\#1 on the final snapshot & 0.9328 & 0.0286 & +65\% & +58\% \\
\bottomrule
\end{tabular}
}

\end{table}

\paragraph{A hypothesis, not a finding.} The final recipe trained at about \val{tokens-per-s}\docnum{} tokens/s on a
\val{params}\docnum{}-parameter model, which we estimated at roughly \val{mfu}\docnum{} model-FLOPs utilisation
\citep{chowdhery2023palm} with about \val{idle}\docnum{} of wall time idle on host dispatch; both are approximate
profiles. At an equal recipe, \cref{eq:exchange} values \val{tput-lo-pct}--\val{tput-hi-pct}\% more tokens per second
at \val{tput-range}~bpb, more than our whole gap to \#10, while late recipe changes were worth about
\val{late-change}\docnum{} each. We therefore \emph{hypothesise} that the cheapest way to close the gap was kernel
throughput: custom NKI kernels \citep{aws_nki_docs} for attention, the fused MLP and the optimizer updates. Nothing in
our data rules out the alternative that better-ranked teams had better recipes or data handling, and the conversion
assumes that every team sits on our compute curve. A test would be to write one kernel, measure its step time with a
short chip screen, price it with \cref{eq:exchange} and confirm it with a cold rehearsal; we did not reach that stage
in Phase~1.
